\section{SHAP (SHapley Additive exPlanations)}

\subsection{Introduction}

\textbf{SHAP (SHapley Additive exPlanations)} was introduced by \citet{Lundberg2017SHAP} as a unified framework to explain individual predictions of any machine learning model.  
It is based on the concept of \textit{Shapley values} from cooperative game theory (\citeauthor{Shapley1953Value}, 1953), which distribute the total “payout” of a game fairly among the players according to their contributions.

In this analogy, the \textit{players} correspond to the input features, and the \textit{payout} corresponds to the model prediction for a given instance.  
The goal of SHAP is to quantify how much each feature contributed to pushing the model’s prediction away from a baseline (the average model output over the dataset).

Although SHAP directly builds on the Shapley framework, it represents a major step forward in usability and adoption:
\begin{itemize}
    \item It provides a connection between \textbf{LIME} and Shapley values via an \textbf{additive feature attribution model}.
    \item It introduces efficient algorithms for specific model classes (e.g., \textbf{TreeSHAP}).
    \item It extends applicability to complex domains such as text and image models (via superpixels or token groups).
    \item It unifies local and global interpretability under a single theoretical foundation.
\end{itemize}

Earlier work by \citet{Strumbelj2011Explaining, Strumbelj2014Explaining} proposed the use of Shapley values for model interpretation, but SHAP popularized the approach by combining a solid theoretical basis with efficient estimation techniques and practical visualization tools.

\subsection{Additive Explanation View}

SHAP represents the explanation as an \textit{additive feature attribution model}:

\begin{equation}
g(z) = \phi_0 + \sum_{i=1}^{M} \phi_i z_i,
\qquad z \in \{0,1\}^M
\end{equation}

Here, $z_i$ indicates whether the feature $i$ is ``present'' ($z_i=1$) or ``absent'' ($z_i=0$) in the coalition.  
The coefficients $\phi_i$ are the \textit{Shapley values}, which represent the contribution of each feature to the prediction.  
This additive structure connects SHAP directly to LIME but ensures that the attributions satisfy the fairness axioms of cooperative game theory.

\subsection{Theoretical Properties}

Shapley values are uniquely defined as the solution satisfying the following properties:

\paragraph{Local Accuracy (Efficiency).}
The sum of feature attributions equals the difference between the prediction and the baseline:

\begin{equation}
\phi_0 = \mathbb{E}[f(X)], 
\qquad
f(x) = \phi_0 + \sum_{i=1}^{M} \phi_i
\end{equation}

\paragraph{Missingness.}
Absent features receive zero attribution:

\begin{equation}
z_i = 0 \Rightarrow \phi_i = 0
\end{equation}

\paragraph{Consistency.}
If a model changes so that a feature’s marginal contribution increases (holding others fixed), its Shapley value cannot decrease:

\begin{equation}
\forall S \subseteq N \setminus \{i\}:\;
f'(x_{S \cup \{i\}}) - f'(x_S) \ge f(x_{S \cup \{i\}}) - f(x_S)
\Rightarrow
\phi_i(f',x) \ge \phi_i(f,x)
\end{equation}

\subsection{Shapley Value Definition}

For a model $f$ and instance $x$, the Shapley value for feature $i$ is defined as:

\begin{equation}
\phi_i(f, x) = 
\sum_{S \subseteq N \setminus \{i\}}
\frac{|S|!\,(M-|S|-1)!}{M!}
\big[v_x(S \cup \{i\}) - v_x(S)\big]
\end{equation}

where $v_x(S)$ is the \textit{value function} of coalition $S$, representing the expected model output when only the features in $S$ are known.

\subsubsection*{Value Functions}

\textbf{Interventional (marginal) form:}
\begin{equation}
v_x(S) = \mathbb{E}_{X_{\bar S}}\!\big[\, f(x_S, X_{\bar S}) \,\big]
\end{equation}

\textbf{Conditional form:}
\begin{equation}
v_x^{\text{cond}}(S) = \mathbb{E}_{X_{\bar S}\mid X_S = x_S}\!\big[\, f(x_S, X_{\bar S}) \,\big]
\end{equation}

The interventional form treats features as independent, while the conditional form accounts for feature dependencies but changes the underlying value function of the game.

\subsection{KernelSHAP}

KernelSHAP is the model-agnostic estimator introduced in the original SHAP paper.  
It approximates Shapley values via a weighted linear regression model over randomly sampled coalitions.

Each coalition is represented as a binary vector $z$ that indicates which features are present.  
A mapping function $h_x(z)$ reconstructs a valid input sample:

\begin{equation}
h_x(z)_i = 
\begin{cases}
x_i, & z_i = 1 \\
\tilde{x}_i, & z_i = 0 \quad (\tilde{x} \text{ sampled from background data})
\end{cases}
\end{equation}

The SHAP kernel assigns weights to each coalition:
\begin{equation}
\pi(z) = \frac{M - 1}{\binom{M}{|z|} \, |z| \, (M - |z|)}
\end{equation}

The Shapley values are obtained by solving the weighted least-squares problem:
\begin{equation}
\mathcal{L}(\phi) =
\sum_{m=1}^{K} \pi(z^{(m)})
\Big[
f\!\big(h_x(z^{(m)})\big)
- \big( \phi_0 + \sum_{i=1}^{M} \phi_i z^{(m)}_i \big)
\Big]^2
\end{equation}

Although KernelSHAP is theoretically appealing and connects SHAP to LIME, it is computationally expensive and assumes feature independence, which can yield unreliable explanations in the presence of correlated features.

\subsection{TreeSHAP}

TreeSHAP (\citeauthor{Lundberg2019TreeSHAP}, 2019) is a fast, exact algorithm for computing SHAP values for tree-based models such as decision trees, random forests, and gradient-boosted trees.  
It reduces the computational complexity from exponential to:

\begin{equation}
\mathcal{O}(T \cdot L \cdot D)
\end{equation}

where $T$ is the number of trees, $L$ the number of leaves, and $D$ the maximum tree depth.

Two formulations exist:
\begin{itemize}
    \item \textbf{Interventional TreeSHAP:} Computes classic Shapley values.
    \item \textbf{Path-dependent TreeSHAP:} Relies on conditional expectations along tree paths (used in earlier implementations).
\end{itemize}

TreeSHAP leverages the tree structure to efficiently compute the marginal contributions of features that affect the prediction, avoiding redundant coalition evaluations.

\subsection{Permutation-Based Estimation}

An alternative sampling-based estimator uses random permutations of features rather than random coalitions.  
For each permutation $\pi$, the marginal contribution of feature $i$ is evaluated as:

\begin{equation}
f(x_{S_\pi(i) \cup \{i\}}) - f(x_{S_\pi(i)})
\end{equation}

The exact Shapley value can be expressed as an average over all $M!$ permutations:

\begin{equation}
\phi_i = \frac{1}{M!} \sum_{\pi \in \Pi}
\Big( f(x_{S_\pi(i) \cup \{i\}}) - f(x_{S_\pi(i)}) \Big)
\end{equation}

In practice, SHAP uses Monte Carlo sampling with forward and backward (antithetic) passes:

\begin{equation}
\hat{\phi}_i = \frac{1}{2K} \sum_{k=1}^{K} 
\Big[
f( x_{S_{\pi_k}(i) \cup \{i\}} ) - f( x_{S_{\pi_k}(i)} )
+ f( x_{S_{\pi_k^{-1}}(i) \cup \{i\}} ) - f( x_{S_{\pi_k^{-1}}(i)} )
\Big]
\end{equation}

This approach significantly reduces computational cost while maintaining efficiency consistency.

\subsection{Global SHAP Interpretations}

Running SHAP for all instances produces an $n \times M$ matrix of Shapley values.  
Aggregating these yields various global interpretability plots.

\paragraph{Feature Importance.}
Global importance is computed as the mean absolute SHAP value:

\begin{equation}
I_j = \frac{1}{n} \sum_{k=1}^{n} \big| \phi_j(x^{(k)}) \big|
\end{equation}

\paragraph{Dependence Plot.}
Each point represents a pair $(x^{(k)}_j, \phi_j(x^{(k)}))$:

\begin{equation}
\big\{ (x^{(k)}_j, \phi_j(x^{(k)})) \big\}_{k=1}^{n}
\end{equation}

This shows both the feature effect and its variance across the dataset.  
Coloring by another feature visualizes feature interactions.

\paragraph{Interaction Values.}
The Shapley interaction index measures the pure pairwise effect between features $i$ and $j$:

\begin{equation}
\phi_{i,j}(f, x) =
\sum_{S \subseteq N \setminus \{i,j\}}
\frac{|S|!\,(M-|S|-2)!}{(M-1)!}
\Big[
v_x(S \cup \{i,j\}) - v_x(S \cup \{i\}) - v_x(S \cup \{j\}) + v_x(S)
\Big]
\end{equation}

\subsection{Strengths}

\begin{itemize}
    \item Solid theoretical foundation grounded in cooperative game theory.
    \item Provides consistent local and global explanations from the same attributions.
    \item Connects Shapley values and LIME under one additive framework.
    \item TreeSHAP enables fast, exact computation for tree models.
    \item Rich visualization suite (force plots, summary plots, dependence plots, interaction heatmaps).
\end{itemize}

\subsection{Limitations}

\begin{itemize}
    \item KernelSHAP is computationally expensive and assumes feature independence.
    \item Marginal SHAP ignores correlations, while conditional SHAP changes the underlying value function.
    \item Path-dependent TreeSHAP can yield unintuitive non-zero attributions for unused features.
    \item Shapley-based explanations can be misinterpreted or manipulated (\citeauthor{Slack2020Fooling}, 2020).
\end{itemize}

\subsection{Software Implementations}

\begin{itemize}
    \item \textbf{Python:} The \texttt{shap} package, integrated with \texttt{scikit-learn}, \texttt{XGBoost}, and \texttt{LightGBM}.
    \item \textbf{R:} Packages \texttt{shapper}, \texttt{fastshap}, and SHAP functionality in \texttt{xgboost}.
    \item \textbf{Advanced interactions:} The \texttt{shapiq} Python package.
\end{itemize}

\subsection{Summary}

SHAP provides a unified, theoretically grounded framework for model interpretability.  
By connecting cooperative game theory with practical estimation and visualization tools, it offers consistent local and global explanations applicable to a wide range of model types.  
Its strength lies in its fairness axioms, efficient tree-based implementation, and ability to unify interpretability methods under one coherent paradigm.


\section{Ceteris Paribus (CP) Plots}

\subsection{Introduction}

\textbf{Ceteris Paribus (CP) plots} \citep{Kuzba2019CeterisParibus} visualize how changes in a single feature affect the prediction for a given data instance, while keeping all other features constant.  
The Latin phrase \textit{ceteris paribus} means “all other things being equal,” reflecting the method’s core principle: vary one feature at a time and observe how the model’s output changes.

Despite their simplicity, CP plots form the foundation of many model-agnostic interpretation techniques.  
They enable an intuitive understanding of a model’s local behavior around a specific instance and serve as building blocks for more advanced methods such as:
\begin{itemize}
    \item \textbf{Individual Conditional Expectation (ICE)} plots — which stack multiple CP curves for all instances.
    \item \textbf{Partial Dependence Plots (PDPs)} — which average CP curves across the dataset.
\end{itemize}

\subsection{Concept and Definition}

Given a model $f$, a feature of interest $x_j$, and a particular observation $x = (x_1, x_2, \dots, x_p)$, the Ceteris Paribus function for $x_j$ is defined as:

\begin{equation}
f^{CP}_x(x_j') = f(x_1, \dots, x_{j-1}, x_j', x_{j+1}, \dots, x_p)
\end{equation}

This function describes how the prediction changes when the $j$-th feature value is replaced by $x_j'$, keeping all other feature values fixed at those of the original observation.  
Plotting $f^{CP}_x(x_j')$ across a grid of $x_j'$ values yields the CP curve.

\subsection{Algorithm}

For a numerical feature $x_j$, the procedure is as follows:

\begin{enumerate}
    \item \textbf{Input:} Model $f$, data point $x$, and feature $x_j$ to vary.
    \item \textbf{Create a grid:} Define an equidistant grid 
    \[
    \{ x_j^{(1)}, x_j^{(2)}, \dots, x_j^{(m)} \}, 
    \quad \text{where typically } \min(x_j) \le x_j^{(1)} < \dots < x_j^{(m)} \le \max(x_j).
    \]
    \item \textbf{Generate modified instances:} For each grid value $x_j^{(k)}$, create a modified instance
    \[
    x^{(k)} = (x_1, \dots, x_{j-1}, x_j^{(k)}, x_{j+1}, \dots, x_p).
    \]
    \item \textbf{Predict:} Compute the corresponding predictions 
    \[
    \hat{y}^{(k)} = f(x^{(k)}).
    \]
    \item \textbf{Visualize:} Plot the pairs $(x_j^{(k)}, \hat{y}^{(k)})$ as a continuous curve.  
    Mark the original feature value $x_j$ with a dot.
\end{enumerate}

For a categorical feature, the algorithm becomes even simpler:

\begin{enumerate}
    \item Obtain the set of unique categories 
    \[
    \mathcal{C}_j = \{ c_1, c_2, \dots, c_r \}.
    \]
    \item For each category $c_\ell$, create a modified instance 
    \[
    x^{(\ell)} = (x_1, \dots, x_{j-1}, c_\ell, x_{j+1}, \dots, x_p),
    \]
    and compute 
    \[
    \hat{y}^{(\ell)} = f(x^{(\ell)}).
    \]
    \item Plot $\hat{y}^{(\ell)}$ versus $c_\ell$ as a bar or dot plot.
\end{enumerate}

The grid size determines the smoothness of the CP curve but also the number of model evaluations.

\subsection{Illustrative Examples}

\paragraph{Example 1: Penguin Classification.}
A random forest model predicts the probability of a penguin being female.  
For an Adelie penguin with bill depth $20.6\,\mathrm{mm}$, a CP plot for bill depth shows that small decreases in bill depth slightly increase $P(\text{female})$, but further decreases sharply reduce it.  
The original feature value is indicated by a dot on the curve.

However, correlations between features (e.g., bill depth, body mass, flipper length) may cause unrealistic combinations when varying one feature while keeping others fixed.  
Thus, caution is advised when interpreting parts of the CP curve far from the original value.

\paragraph{Example 2: Bike Rental Prediction.}
An SVM model predicts the number of rented bikes based on weather and seasonal variables.  
CP plots across all features reveal that the number of bikes rented two days earlier and temperature have the strongest influence on predictions.  
For unrealistic manipulations (e.g., increasing temperature to $30^\circ$C during winter), CP plots can yield unrealistic predictions, reflecting model behavior outside the data manifold.

\subsection{Relationship to ICE and PDP}

CP plots are the foundation for more complex model-agnostic visualization methods:

\begin{itemize}
    \item \textbf{ICE Plots:} Contain all CP curves for every observation in the dataset.  
    Formally, an ICE plot is 
    \[
    \{ f^{CP}_{x^{(i)}}(x_j') \}_{i=1}^n.
    \]
    \item \textbf{PDPs:} Represent the average of all CP curves:
    \[
    \text{PDP}(x_j') = \frac{1}{n} \sum_{i=1}^n f^{CP}_{x^{(i)}}(x_j').
    \]
\end{itemize}

Hence, CP plots visualize the local prediction surface for a single instance, ICE plots visualize individual-level heterogeneity, and PDPs provide global trends.

\subsection{Applications and Extensions}

\begin{itemize}
    \item \textbf{Model comparison:} Overlay CP curves from different models (e.g., linear model, SVM, decision tree) to study how each model responds to changes in a feature.
    \item \textbf{Feature comparison:} Plot CP curves for multiple features to identify which features influence the prediction most strongly.
    \item \textbf{Class comparison:} For multi-class models, plot separate CP curves for each class probability.
    \item \textbf{Data subsets:} Compute CP curves on subgroups (e.g., by gender, region) to study differential model behavior.
\end{itemize}

\subsection{Interpretation Considerations}

\begin{itemize}
    \item \textbf{Feature dependence:}  
    When features are correlated, modifying one feature while keeping others fixed can create unrealistic data points.  
    Analysts should restrict the range of $x_j'$ to realistic intervals or use conditional sampling.
    \item \textbf{Causal interpretation:}  
    CP plots describe correlations in the model’s learned function, not necessarily causal effects.  
    Causal conclusions require a causal model structure or experimental design.
\end{itemize}

\subsection{Visualization Variants}

\begin{itemize}
    \item \textbf{Full CP curves:} Continuous lines showing prediction changes across the entire feature range.
    \item \textbf{Sparklines:} Minimalist line plots summarizing local prediction sensitivity (popularized by \citet{Tufte1983Sparklines}).
\end{itemize}

\subsection{Strengths}

\begin{itemize}
    \item \textbf{Simplicity:}  
    Easy to compute and interpret; ideal for introducing model-agnostic explainability.
    \item \textbf{Complementarity:}  
    Complements attribution methods (e.g., SHAP, LIME) by showing prediction sensitivity to feature changes.
    \item \textbf{Flexibility:}  
    Serves as a foundational building block for ICE and PDP plots and can be extended across models, hyperparameters, and classes.
\end{itemize}

\subsection{Limitations}

\begin{itemize}
    \item \textbf{One-feature-at-a-time:}  
    CP plots only vary one feature, ignoring interactions unless manually extended.
    \item \textbf{Feature correlation:}  
    Unrealistic scenarios can arise when correlated features are treated as independent.
    \item \textbf{Risk of misinterpretation:}  
    Non-experts may incorrectly interpret CP curves causally; careful communication is necessary.
\end{itemize}

\subsection{Summary}

Ceteris Paribus plots are simple yet powerful tools for understanding how model predictions respond to feature changes.  
They provide clear visual explanations for individual instances, support model comparison, and form the foundation of more complex interpretability methods like ICE and PDP.  
Despite their simplicity, they should be applied carefully in the presence of correlated features or when causal interpretations are implied.


